
% Kapitola 2: Metadáta vedeckých výstupov a authority control
% Stav: finálna verzia

\section{Metadáta vedeckých výstupov a authority control}

Pre návrh aplikácie na evidenciu a normalizáciu publikačných záznamov je potrebné najprv presne vymedziť, aké typy údajov sú predmetom spracovania a akú úlohu plnia v prostredí inštitucionálneho repozitára. Praktická časť nebude pracovať iba s voľným textom, ale s čiastočne štruktúrovanými bibliografickými údajmi rôznej kvality, ktoré sa budú porovnávať, dopĺňať, validovať a prepájať s lokálnymi autoritami. Táto kapitola preto zavádza pojmy a referenčné rámce, bez ktorých by nebolo možné korektne opísať deduplikáciu, normalizáciu ani auditovateľné spracovanie záznamov.

\subsection{Pojem metadát a ich rola v repozitári}

Metadáta možno v najvšeobecnejšom zmysle chápať ako štruktúrované údaje o zdroji. V prostredí vedeckých výstupov ide predovšetkým o údaje, ktoré umožňujú dokument identifikovať, opísať, vyhľadať, previazať s ďalšími entitami a sprístupniť v prepojiteľnom prostredí \cite{DCMITerms}. Pre inštitucionálny repozitár preto metadáta nepredstavujú iba sprievodný opis dokumentu, ale základnú vrstvu, na ktorej stojí evidencia, vyhľadávanie, deduplikácia aj výmena údajov s externými systémami.

V tejto práci sú dôležité najmä tri roviny metadát. Prvú tvoria deskriptívne údaje, ako názov, autori, názov časopisu, rok vydania alebo identifikátory. Druhú rovinu predstavujú väzbové a autoritné údaje, ktoré prepájajú textové polia s konkrétnymi osobami, periodikami alebo inými kontrolovanými entitami \cite{DCMITerms}. Tretiu rovinu tvoria prevádzkové údaje potrebné pre interné spracovanie, napríklad pôvod hodnoty, stav validácie alebo informácia o tom, či bola hodnota doplnená heuristikou, modelom alebo používateľom. Táto tretia rovina je vymedzená pre potreby tejto práce a bude podrobnejšie rozpracovaná v podkapitole venovanej provenancii a auditovateľnosti úprav.

\subsection{Bibliografické metadátové schémy}

Pri spracovaní vedeckých výstupov sa nepoužíva jediná univerzálna schéma. V praxi sa stretávajú všeobecné metadátové modely, knižnične orientované schémy aj prenosové formáty určené na export a výmenu medzi nástrojmi. Pre účely tejto práce je dôležité rozlišovať najmä medzi tým, čo predstavuje konceptuálnu schému opisu, a tým, čo predstavuje technický formát výmeny.

Dublin Core patrí medzi najrozšírenejšie všeobecné metadátové rámce. Jeho význam spočíva najmä v tom, že poskytuje relatívne jednoduchý a prenositeľný súbor prvkov použiteľných naprieč rôznymi typmi repozitárov a informačných služieb \cite{DCMITerms}. Pre bohatší bibliografický opis je však často príliš všeobecný. V tomto smere je podrobnejšia schéma MODS, ktorá vznikla v knižničnom prostredí a umožňuje presnejšie zachytiť bibliografické a vzťahové charakteristiky dokumentu \cite{MODSUserGuide}. Rozdiel medzi týmito prístupmi je pre túto prácu dôležitý preto, že importované záznamy môžu obsahovať údaje vhodné na jednoduchú výmenu, no zároveň vyžadujú detailnejšiu internú reprezentáciu pre potreby kurácie a normalizácie.

Popri metadátových schémach treba rozlišovať aj exportné a citačné formáty. BibTeX je formát pevne spätý s prostredím \LaTeX{} a slúži predovšetkým na uchovávanie a generovanie bibliografických záznamov v citačných workflow \cite{Patashnik1988}. V bibliografických databázach sa však používa aj formát RIS, ktorý v praxi slúži ako rozšírený výmenný formát medzi databázami, citačnými manažérmi a ďalšími nástrojmi. Pre túto prácu nie je rozhodujúca jeho úplná syntaktická špecifikácia, ale skutočnosť, že rôzne vstupné formáty vyjadrujú podobné bibliografické údaje odlišným spôsobom a že aplikácia bude musieť tieto rozdiely preklenúť pri importe a transformácii.

\subsection{Perzistentné identifikátory}

Osobitné miesto v správe metadát majú perzistentné identifikátory. Ich hlavnou úlohou je zabezpečiť, aby bolo možné konkrétnu entitu jednoznačne a dlhodobo identifikovať naprieč systémami, nezávisle od toho, v akej textovej podobe je opísaná v konkrétnom zázname. V rámci riešenej domény ide najmä o identifikátory dokumentov, autorov a seriálových zdrojov.

Pri jednotlivých publikáciách zohráva kľúčovú úlohu DOI. Tento identifikátor je v praxi najspoľahlivejším oporným bodom pri deduplikácii záznamov, pretože je navrhnutý práve na jednoznačnú identifikáciu objektu a jeho strojové rozlíšenie \cite{DOIFoundation,ISO26324_2025}. V prostredí evidencie publikácií preto predstavuje prirodzený primárny kľúč, ak je v zázname dostupný a syntakticky korektný. Pri autoroch má analogickú, hoci nie úplne univerzálnu, úlohu ORCID. Ten nepokrýva všetkých autorov, ale tam, kde je k dispozícii, výrazne znižuje nejednoznačnosť pri párovaní mien a pri prepájaní záznamov medzi systémami \cite{ORCIDWhatIs}. V prípade periodík slúžia ako stabilizačný prvok ISSN a e-ISSN, ktoré umožňujú odlíšiť samotný časopis od variantných textových zápisov jeho názvu \cite{ISSNWhatIs,ISO3297_2022}.

Z pohľadu praktickej časti je dôležité, že identifikátory nevyriešia celú úlohu samy osebe. Nie všetky importované záznamy ich obsahujú, nie vždy sú uvedené v správnom formáte a nie všetky polia sú medzi zdrojmi konzistentné. Napriek tomu práve identifikátory tvoria základný referenčný rámec, voči ktorému sa budú posudzovať variantné zápisy v textových poliach. Ich prítomnosť preto zásadne ovplyvňuje, či sa záznam bude dať riešiť deterministicky, alebo bude potrebné prejsť na porovnanie podobnosti.

\subsection{Authority control}

Authority control predstavuje súbor princípov a postupov, ktorých cieľom je zabezpečiť konzistentné reprezentovanie entít v bibliografickom a repozitárovom prostredí. Nejde iba o technický register mien, ale o spôsob, ako oddeliť samotnú entitu od jej textových variantov a ako nad touto entitou dlhodobo udržiavať kontrolovaný opis. V modeli FRAD je authority data chápané ako základ pre riadené prístupové body a pre vzťahy medzi rôznymi podobami mien, názvov a identít \cite{FRAD2013}.

V kontexte tejto práce sa authority control týka predovšetkým dvoch oblastí. Prvou sú autori, pri ktorých môže mať tá istá osoba viacero zápisov mena, rôzne afiliácie alebo iba čiastočne uvedené iniciály. Druhou sú časopisy, ktorých názvy sa v exportoch môžu objavovať v plnej aj skrátenej podobe, s interpunkčnými odchýlkami alebo s rozdielnym uvedením identifikátorov. Pre praktické riešenie je preto potrebné vychádzať z lokálnych autoritných záznamov a nevnímať textové pole ako definitívny opis entity. Externé služby, ako ORCID, ISSN Portal, Crossref alebo VIAF, môžu slúžiť ako doplnkový referenčný zdroj, no finálna reprezentácia musí byť v lokálnom workflow kontrolovaná vzhľadom na potreby UTB \cite{ORCIDWhatIs,ISSNWhatIs,VIAF,CrossrefRESTAPI}.

Z hľadiska návrhu aplikácie je authority control dôležitý aj preto, že vytvára most medzi normalizáciou textu a evidenčnou integritou systému. Cieľom nie je iba prepísať nejednotný reťazec do esteticky lepšej podoby, ale priradiť ho ku konkrétnej entite, ktorú možno ďalej jednoznačne používať pri vyhľadávaní, validácii, deduplikácii a reportovaní. Táto logika bude v praktickej časti priamo prítomná pri mapovaní autorov na interný register UTB autorov a pri normalizácii časopisov.

\subsection{Externé zdroje pravdy}

Pri správe publikačných záznamov nemožno predpokladať, že jeden zdroj poskytne vždy úplné a bezchybné metadáta. Preto je vhodnejšie hovoriť o externých zdrojoch pravdy v pluráli a chápať ich ako referenčné body s rôznou mierou autority pre rôzne typy údajov. V podmienkach UTB majú osobitné postavenie Web of Science a Scopus, pretože predstavujú hlavné vstupné zdroje importu. Ich výhodou je rozsiahle pokrytie a štandardizované možnosti exportu, nevýhodou však je, že výsledné záznamy nemusia zodpovedať lokálnym požiadavkám na autority, interné väzby ani finálnu podobu evidencie \cite{ClarivateWOSExport,ScopusSupport}.

Významnú doplnkovú rolu má Crossref. Jeho rozhranie REST API umožňuje vyhľadávať a získavať bibliografické metadáta k DOI, periodikám aj jednotlivým dielam vo forme vhodnej na ďalšie strojové spracovanie \cite{CrossrefRESTAPI}. Crossref zároveň explicitne uvádza, že jeho metadáta sú voľne dostupné pre komunitu a že väčšina bibliografických metadát je považovaná za fakty nepodliehajúce autorskému právu, pričom špecifické obmedzenia sa týkajú najmä abstraktov \cite{CrossrefMetadataRetrieval}. Pre riešenú aplikáciu je preto Crossref vhodným doplnkovým zdrojom najmä tam, kde treba overiť alebo doplniť časopis, DOI alebo ďalšie chýbajúce bibliografické polia.

Otvorený charakter má aj OpenAlex, ktorý poskytuje verejne dostupné rozhranie nad rozsiahlym agregovaným grafom vedeckých entít \cite{OpenAlexDocs}. Jeho odborné zarámcovanie poskytuje aj práca Priema, Piwowar a Orra, ktorá OpenAlex opisuje ako plne otvorený index vedeckých diel, autorov, časopisov, inštitúcií a konceptov \cite{OpenAlexPaper}. Z pohľadu tejto práce však OpenAlex nepredstavuje primárny importný kanál, ale skôr pomocný zdroj na overovanie alebo dopĺňanie vybraných údajov. Praktická časť preto nebude vychádzať z predstavy jediného absolútneho zdroja pravdy, ale z kombinácie primárnych importných databáz, otvorených referenčných služieb a lokálnych autoritných záznamov. Práve toto viaczdrojové poňatie je predpokladom pre spoľahlivú normalizáciu metadát v prostredí univerzitnej knižnice.
